\begin{figure}[htbp]
\centering
\begin{tikzpicture}[x=1.55cm, y=0.72cm]

\tikzset{
  done/.style    = {fill=VIUOrange!70, draw=VIUDeep,    line width=0.45pt, rounded corners=2pt},
  current/.style = {fill=VIUOrange,    draw=VIUDeep,    line width=0.75pt, rounded corners=2pt},
  pending/.style = {fill=VIUGray!12,   draw=VIUGray!45, line width=0.45pt, rounded corners=2pt},
  defense/.style = {fill=VIUDeep!18,   draw=VIUDeep,    line width=0.75pt, rounded corners=2pt}
}
\def\BH{0.50}

%-- Franjas alternadas de fondo (cada dos meses para guiar la lectura)
\foreach \m in {0,2,4,6}{
  \fill[VIUGray!7] (\m, 0.68) rectangle (\m+1, 7.58);
}
%-- Separadores de mes
\foreach \m in {0,...,7}{
  \draw[VIUGray!30, very thin] (\m, 0.68) -- (\m, 7.58);
}
%-- Línea de base horizontal
\draw[VIUGray!35, thin] (0, 0.68) -- (7, 0.68);

%-- Etiquetas de meses (encabezado superior)
\foreach \m/\lbl in {0/Mar,1/Abr,2/May,3/Jun,4/Jul,5/Ago,6/Sep}{
  \node[font=\footnotesize\bfseries, color=VIUGray!80!black, anchor=south]
    at (\m+0.5, 7.60) {\lbl};
}
\node[font=\small\bfseries, color=VIUDark, anchor=south] at (3.5, 8.10) {2026};

%-------------------------------------------------------------
%  BARRAS DE FASES + etiquetas (derecha de la barra, izquierda del diagrama)
%-------------------------------------------------------------

% E1 — Revisión bibliográfica  [Mar 4 → May]  COMPLETADO
\fill[done] (0.10, 7.0) rectangle (2.0, 7.0+\BH);
\node[font=\footnotesize, color=VIUDark, anchor=east, text width=3.7cm, align=right]
  at (-0.10, 7.25) {Revisión bibliográfica};

% E2 — Formulación del problema  [Mar 4 → 29 may]  ENTREGADO (esta Tarea)
\fill[current] (0.10, 6.0) rectangle (2.9, 6.0+\BH);
\node[font=\footnotesize, color=VIUDark, anchor=east, text width=3.7cm, align=right]
  at (-0.10, 6.25) {Formulación del problema};

% E3 — Implementación Python  [Jun → mid-ago]  PENDIENTE
\fill[pending] (2.0, 5.0) rectangle (4.5, 5.0+\BH);
\node[font=\footnotesize, color=VIUDark, anchor=east, text width=3.7cm, align=right]
  at (-0.10, 5.25) {Implementación Python};

% E4 — Simulación experimental  [Jun → ago]  PENDIENTE
\fill[pending] (3.1, 4.0) rectangle (5.5, 4.0+\BH);
\node[font=\footnotesize, color=VIUDark, anchor=east, text width=3.7cm, align=right]
  at (-0.10, 4.25) {Simulación experimental};

% E5 — Validación en CoppeliaSim  [mid-ago → sep]  PENDIENTE
\fill[pending] (4.5, 3.0) rectangle (6.0, 3.0+\BH);
\node[font=\footnotesize, color=VIUDark, anchor=east, text width=3.7cm, align=right]
  at (-0.10, 3.25) {Validación CoppeliaSim};

% E6 — Redacción de la memoria  [jul → sep]  PENDIENTE
\fill[pending] (4.0, 2.0) rectangle (6.7, 2.0+\BH);
\node[font=\footnotesize, color=VIUDark, anchor=east, text width=3.7cm, align=right]
  at (-0.10, 2.25) {Redacción de la memoria};

% E7 — Defensa TFM  [21–25 sep]  PENDIENTE
\fill[defense] (6.5, 1.0) rectangle (7.0, 1.0+\BH);
\node[font=\footnotesize, color=VIUDark, anchor=east, text width=3.7cm, align=right]
  at (-0.10, 1.25) {Defensa TFM};

%-------------------------------------------------------------
%  TUTORÍAS COLECTIVAS  (triángulos sobre el límite superior)
%  T0 = 4 mar (x≈0.10)   T1 = 3 jun (x≈3.08)   T2 = 15 jul (x≈4.47)
%-------------------------------------------------------------
\foreach \tx in {0.10, 3.08, 4.47}{
  \fill[VIUGray!55]
    (\tx, 7.58) -- (\tx-0.08, 7.74) -- (\tx+0.08, 7.74) -- cycle;
}
\node[font=\tiny\bfseries, color=VIUGray!80, anchor=south west] at (0.18,  7.74) {T0};
\node[font=\tiny\bfseries, color=VIUGray!80, anchor=south west] at (3.16,  7.74) {T1};
\node[font=\tiny\bfseries, color=VIUGray!80, anchor=south west] at (4.55,  7.74) {T2};

%-------------------------------------------------------------
%  HITOS
%-------------------------------------------------------------

% Tarea 1 / Informe Intermedio  (29 may → x = 2.9)
\draw[VIUDeep, line width=1.5pt] (2.9, 0.68) -- (2.9, 7.58);
\fill[VIUDeep] (2.9, 0.68) circle(3.2pt);
\node[font=\tiny\bfseries, color=VIUDeep, anchor=north, align=center]
  at (2.9, 0.62) {Tarea\,1\\29~may.};

% Defensa, 1\textordfeminine{} convocatoria  (21 sep → x = 6.67)
\draw[VIUDark!55, line width=1.0pt, dashed] (6.67, 0.68) -- (6.67, 7.58);
\fill[VIUDark!55] (6.67, 0.68) circle(2.5pt);
\node[font=\tiny, color=VIUDark!70, anchor=north east, align=right]
  at (6.67, 0.62) {Defensa\\21--25~sep.};

%-------------------------------------------------------------
%  LEYENDA
%-------------------------------------------------------------
\fill[done]    (0.00,-0.72) rectangle (0.90,-0.35);
\node[font=\footnotesize, color=VIUDark, anchor=west] at (1.00,-0.535) {Completado};
\fill[current] (2.60,-0.72) rectangle (3.50,-0.35);
\node[font=\footnotesize, color=VIUDark, anchor=west] at (3.60,-0.535) {Entregado};
\fill[pending] (5.10,-0.72) rectangle (6.00,-0.35);
\node[font=\footnotesize, color=VIUDark, anchor=west] at (6.10,-0.535) {Pendiente};

\end{tikzpicture}
\caption[Cronograma del TFM]%
{Cronograma del TFM (marzo--septiembre 2026).
Tutorías colectivas síncronas: T0~inicio (4~mar.), T1~intermedia (3~jun.),
T2~final (15~jul.).
Fases: E1~Revisión bibliográfica, E2~Formulación del problema (Tarea~1),
E3~Implementación Python, E4~Simulación experimental,
E5~Validación CoppeliaSim, E6~Redacción de la memoria,
E7~Defensa (1.\textordfeminine\ conv., 21--25~sep.\ 2026).
La línea sólida marca el hito actual (29~may.\ 2026); la de trazos,
la convocatoria de defensa.}
\label{fig:gantt}
\end{figure}
